\documentclass[UTF8,fontset=windows,11pt]{ctexart}
\usepackage[a4paper,margin=2.35cm]{geometry}
\usepackage{amsmath,amssymb,bm,mathtools}
\usepackage{booktabs,longtable,array}
\usepackage{graphicx}
\graphicspath{{../coding/effective-social-input/results/paper-kf-extended-v3/full/}%
{../coding/effective-social-input/results/paper-simulation-figures/}%
{../coding/effective-social-input/results/mainstream-model-comparison/}%
{../coding/effective-social-input/results/turn-occlusion-pilot-v1/full/}%
{../coding/effective-social-input/results/turn-occlusion-pilot-v1/smoke/}}
\usepackage{xcolor}
\usepackage{enumitem}
\usepackage{algorithm}
\usepackage{algpseudocode}
\usepackage{amsthm}
\usepackage{caption}
\usepackage{hyperref}
\usepackage{microtype}
\newcommand{\upcite}[1]{\textsuperscript{\cite{#1}}}
\floatname{algorithm}{算法}
\newtheorem{theorem}{定理}[section]
\newtheorem{proposition}[theorem]{命题}
\newtheorem{lemma}[theorem]{引理}
\newtheorem{corollary}[theorem]{推论}
\newtheorem{definition}[theorem]{定义}
\newcommand{\Sone}{\mathbb S^1}
\newcommand{\Ttwo}{\mathbb T^2}
\newcommand{\E}{\mathbb E}
\newcommand{\R}{\mathbb R}
\algrenewcommand\algorithmicrequire{\textbf{输入：}}
\algrenewcommand\algorithmicensure{\textbf{输出：}}
\algrenewcommand\algorithmicfor{\textbf{对于}}
\algrenewcommand\algorithmicdo{\textbf{执行}}
\algrenewcommand\algorithmicend{\textbf{结束}}
\algrenewcommand\algorithmicreturn{\textbf{返回}}
\hypersetup{colorlinks=true,linkcolor=black,citecolor=black,urlcolor=blue,
 pdftitle={间歇方位观测下异中心神经场集群运动的Kalman预测补偿方法},
 pdfsubject={异中心神经场集群运动与Kalman预测补偿},
 pdfkeywords={异中心表征；环形吸引子；集群运动；间歇观测；Kalman滤波},
 bookmarksnumbered=true}
\setlength{\parindent}{2em}
\setlength{\parskip}{0.25em}
\linespread{1.18}
\setlength{\emergencystretch}{1em}
\title{\bfseries 间歇方位观测下异中心神经场集群运动的\newline Kalman预测补偿方法}
\author{}
\date{}
\begin{document}
\maketitle

\begin{abstract}
异中心（allocentric）环形吸引子模型能够仅依靠邻居方位输入产生集群有序运动，但其性能依赖连续且可靠的社会感知。当局部遮挡或通信中断导致邻居方位观测间歇缺失时，神经场输入会减少或发生方向偏差，进而引起活动峰扩散、运动方向失序和集群结构退化。本文首先在 Salahshour 与 Couzin 的异中心神经场 flocking 模型基础上，分析社会输入数量、输入幅值和方向覆盖对集群运动的影响；随后针对当前模型仅具有方位观测的特点，建立方位—方位角速度二状态 Kalman 预测器，在观测缺失期间外推已跟踪邻居的方位，并将预测协方差通过周期核卷积传递至神经场输入。本文推导了连续缺失期间方位协方差的增长规律、圆周方向信息的衰减规律以及条件均方输入重建性质，明确了该方法的适用条件和局限。实验比较新鲜观测删除、最后方位保持、Kalman均值预测和协方差感知预测，并使用全新种子的配对闭环实验评估遮挡强度和持续时间的影响。为避免在已经收敛的直线运动状态下删除冗余邻居而低估预测价值，进一步设置遮挡前受控转向实验，并以相对无遮挡控制的退化量作为主响应。结果表明，在固定身份、等权、原始单边增益条件下，Kalman补偿能够降低连续缺失造成的相对性能偏差；阻尼角速度KF还可进一步降低隐藏边的方位预测误差。
\end{abstract}

\noindent\textbf{关键词：} 异中心表征；环形吸引子；集群运动；间歇观测；Kalman滤波；方位预测；不确定性传播

\section{引言}

无人集群通过多个自主个体的局部感知、相互协同和分布式决策，能够完成单体平台难以承担的大范围、长时程和高风险任务。在灾害监测与应急搜救中，多无人机集群可在有限视场和避障约束下完成协同覆盖与目标搜寻\upcite{Bai2024}；在林下环境巡检与农业作业中，分布式集群控制能够提高复杂环境中的队形保持和环境适应能力\upcite{Boldrer2026}；在通信中继等场景中，个体通信障碍会直接影响集群凝聚和任务执行\upcite{LiFlocking2024}，而间歇采样和通信条件下的协调控制则是维持区域覆盖与群体协同的关键\upcite{Ge2025Coordination}。因此，无人集群在上述领域具有重要应用价值。

若按决策所依赖的信息和计算组织方式划分，多智能体协同控制可分为集中式与分布式两类。集中式方法汇集个体状态并统一计算控制指令，便于协调全局任务，但需要持续获取和传递群体信息；分布式方法则由各个体依据本地观测及可获得的邻居信息更新自身行为\upcite{OlfatiSaber2006}。对于需要持续协同运动的无人集群，后者能够将群体协调建立在局部交互之上，但其运动效果也取决于邻居信息的可获得性和表达方式。由局部交互产生有序集体运动的 flocking，因而成为研究这一问题的典型模型。

经典 flocking 模型通常直接规定个体如何根据邻居运动：Vicsek 模型通过局部航向对齐研究自驱动粒子的有序相变\upcite{Vicsek1995}，Olfati--Saber 则将邻居间的相互作用写成分布式控制律\upcite{OlfatiSaber2006}。这类模型描述了局部信息如何影响运动，却较少讨论邻居方向信息在个体内部采用何种参照系表征。Salahshour 与 Couzin 从这一表征问题出发，将邻居方位作为环形吸引子网络的感觉输入，比较相对于自身朝向的自我中心表征与相对于外部稳定参照的异中心表征；在其模型中，异中心表征无需预设显式对齐规则即可产生有序的集体运动\upcite{Salahshour2025}。因此，异中心神经场为研究仅依靠局部方位输入形成集群运动提供了具体机制。

然而，原始异中心神经场模型以邻居方位持续可得为前提，遮挡或链路中断会使这一社会输入间歇缺失。直接删除邻居输入会同时改变输入幅值与方位覆盖；沿用最后一次观测虽能保留输入幅值，却可能在群体转向时持续提供过期方向。两种处理对神经场及群体运动的影响并不等价，也不能仅由当前可见邻居数解释。更重要的是，仅有方位观测时，个体无法在缺失期间直接获得邻居的距离或二维位置。因此，这里的预测对象应是神经场所需的方向输入，而非不可由方位直接恢复的邻居位置。

本文关注的核心问题是：在间歇仅方位观测下，如何利用历史方位延续异中心神经场的有效社会输入，并在有限缺失窗口内减轻群体运动相对无遮挡状态的退化？实现这一目标需要解决三个递进的科学问题。其一，辨明邻居数量、总输入幅值和方位结构各自对群体有序性与凝聚性的作用，识别观测缺失造成的信息损失。其二，在已知自身朝向且邻居关联正确的前提下，确定历史方位可支持的因果预测形式，并建立预测不确定性到周期社会输入的映射。其三，厘清输入重建的局部效果向闭环群体运动传递的条件，以及相对最后方位保持的增量与失效边界。

为回答第一个问题，本文以原始异中心神经场为基线，考察邻居数量、总输入幅值和方位覆盖与群体有序度、凝聚程度的关系，并比较删除输入和保持最后方位造成的差异。

在此基础上，针对第二个问题，本文对已关联邻居建立方位--方位角速度 Kalman 预测器，在观测缺失时传播预测状态，将方位均值和协方差映射为圆周方向输入，并分析方位方差增长、方向信息衰减及输入重建误差的条件性界。

针对第三个问题，本文采用共同检查点分叉的配对实验，在不同连续缺失时长和受控转向条件下，将预测方法与删除输入、保持旧方位进行比较，分别评估方位误差、全局与局部有序度及凝聚偏差，界定均值外推和不确定性传播的条件性作用。上述研究以已关联邻居和给定遮挡过程为前提，其结论限于相应的模型与实验条件。
\clearpage

\section{问题描述}

\subsection{问题描述}

考虑 $N$ 个智能体在边长 $L$ 的二维周期区域内运动。每个智能体 $i$ 具有物理状态
\begin{equation}
 (q_i(k),\phi_i(k)),
\end{equation}
其中 $q_i(k)=[x_i(k),y_i(k)]^\top$ 是位置，$\phi_i(k)$ 是世界坐标系中的自身朝向。每个智能体同时维护一个包含 $M$ 个方向单元的环形神经场状态 $u_i(k)\in\mathbb R^M$ 和固定方向偏好
\begin{equation}
 \alpha_m=\frac{2\pi m}{M}+\phi_{i,0},\qquad m=0,\ldots,M-1,
\end{equation}
其中 $\phi_{i,0}$ 是初始化方向偏移。该表示是异中心的，即神经元方向固定在共同世界坐标系中，不随智能体自身朝向旋转。

智能体不能直接获得邻居位置，也不具有距离传感器。对于每条有向邻居边 $(i,j)$，传感器可能输出机体坐标系中的方位
\begin{equation}
 z^b_{ij}(k)=
 \operatorname{wrap}\left(\theta_{ij}(k)-\phi_i(k)+\nu_{ij}(k)\right),
\end{equation}
其中 $\theta_{ij}(k)$ 是邻居的世界方位，$\nu_{ij}(k)\sim\mathcal N(0,R)$ 是测量噪声。定义观测指示量
\begin{equation}
 \delta_{ij}(k)=
 \begin{cases}
 1, & \text{第 $k$ 步获得邻居 $j$ 的新鲜方位},\\
 0, & \text{观测缺失或通信中断}.
 \end{cases}
\end{equation}
已知自身朝向时，可将测量转换到世界坐标系：
\begin{equation}
 z_{ij}(k)=\operatorname{wrap}(z^b_{ij}(k)+\phi_i(k)).
\end{equation}

本文的问题可以描述为：在只具有间歇方位观测的条件下\upcite{Xiourouppa2025}，构造一个因果预测与输入重建接口，使神经场在观测缺失期间仍能获得尽可能接近原始社会输入的方向信息，从而维持群体有序运动、凝聚结构和转向跟随能力，而控制器不得读取邻居的真实位置、真实方位或距离。

该问题的输入和输出如图~\ref{tab:problem_io} 所示。

\begin{table}[htbp]
 \centering
 \small
 \caption{问题描述中的控制器输入与输出}
 \label{tab:problem_io}
 \begin{tabular}{>{\raggedright\arraybackslash}p{0.16\linewidth}
                 >{\raggedright\arraybackslash}p{0.37\linewidth}
                 >{\raggedright\arraybackslash}p{0.35\linewidth}}
 \toprule
 层级 & 输入 & 输出\\
 \midrule
 传感器层 &
 自身朝向 $\phi_i(k)$；已关联邻居编号 $j$；新鲜方位 $z_{ij}(k)$；观测指示量 $\delta_{ij}(k)$ &
 经过噪声和遮挡处理后的方位观测\\
 预测层 &
 当前边历史观测 $z_{ij}(0:k)$；上一步估计 $\hat x_{ij}(k-1)$ 和协方差 $P_{ij}(k-1)$ &
 预测方位均值 $\mu_{ij}(k)$；方位方差 $V_{ij}(k)$；轨迹有效指示量 $m_{ij}(k)$\\
 神经场层 &
 神经元方向 $\alpha_m$；神经场状态 $u_i(k)$；预测方位及方差 &
 外部社会输入 $I_i(k)$；更新后的神经活动 $u_i(k+1)$\\
 运动层 &
 正神经活动 $a_{im}(k)$；速度尺度 $v_s$ &
 运动速度 $v_i(k)$；更新位置 $q_i(k+1)$\\
 \bottomrule
 \end{tabular}
\end{table}

令 $\mathcal Z_i(0:k)=\{(j,z_{ij}(t),\delta_{ij}(t)):t\le k\}$ 表示智能体 $i$ 截至第 $k$ 步的全部可用观测历史。控制器的目标可以写为
\begin{equation}
 \min_{\Pi}\ \mathbb E\left[
 \left\|I_i^{\mathrm{true}}(k)-\Pi(\mathcal Z_i(0:k),u_i(k);\Theta)\right\|_2^2
 \right],
\end{equation}
其中 $\Pi$ 是满足因果性约束的预测输入重建接口，$\Theta$ 是滤波器参数，$I_i^{\mathrm{true}}$ 是如果所有邻居方位均新鲜时原模型将获得的社会输入。除输入重建误差外，闭环性能还由 GO、LO、凝聚距离、转向跟随误差和恢复时间评价。

本文进一步假设：

\begin{enumerate}[label=(\arabic*)]
 \item 邻居编号已经正确关联，不考虑错误匹配和轨迹交换；
 \item 自身朝向 $\phi_i(k)$ 已知，不考虑自身朝向估计误差；
 \item 传感器只提供方位，不提供距离，因此不允许使用二维位置—速度滤波器；
 \item 预测接口只能使用截至当前时刻的历史观测，不能使用仿真环境中的未来状态或隐藏真值。
\end{enumerate}

\subsection{原始异中心神经场模型的求解}

原始模型假定邻居社会输入连续可得，并使用全邻居交互。对每个智能体 $i$，其社会邻居集合固定为
\begin{equation}
 \mathcal B_i=\{j\in\{1,\ldots,N\}:j\ne i\}.
\end{equation}
对邻居 $j$，采用周期区域中的最小镜像差分
\begin{equation}
 r_{ij}(k)=q_j(k)-q_i(k)-L\operatorname{round}\left(\frac{q_j(k)-q_i(k)}{L}\right).
\end{equation}
其真实世界方位为
\begin{equation}
 \theta_{ij}(k)=\operatorname{atan2}(r_{ij,y}(k),r_{ij,x}(k)).
\end{equation}
定义圆周角距离
\begin{equation}
 d_{\mathbb S^1}(a,b)=\left|\operatorname{wrap}(a-b)\right|,
\quad \operatorname{wrap}(x)=\operatorname{mod}(x+\pi,2\pi)-\pi.
\end{equation}

原始社会输入为
\begin{equation}
 I^{\mathrm{soc}}_{im}(k)=h_0\sum_{j\in\mathcal B_i}
 g_\sigma\left(\alpha_m-\theta_{ij}(k)\right),
 \quad h_0=\frac{h_t}{N},
\end{equation}
其中
\begin{equation}
 g_\sigma(x)=\exp\left[-\frac{d_{\mathbb S^1}^2(x,0)}{2\sigma_s^2}\right].
\end{equation}
这里 $\mathcal B_i$ 是固定全邻居集合。实际幅值总和为 $(N-1)h_t/N$。

环形连接矩阵与神经场更新分别为
\begin{equation}
 W_{mn}=\cos\left[\pi\left(\frac{d_{mn}}{\pi}\right)^\zeta\right],
\end{equation}
\begin{equation}
 u_i(k+1)=u_i(k)+\Delta t\left[-u_i(k)+\frac{1}{M}W\tanh(\beta u_i(k))-h_b\bm 1+I_i(k)\right].
\end{equation}
神经场内部使用未截断的 $\tanh$。运动解码时定义正活动
\begin{equation}
 a_{im}(k)=\max\{\tanh(\beta u_{im}(k)),0\},
\end{equation}
并令
\begin{equation}
 c_i(k)=\sum_{m=0}^{M-1}a_{im}(k)e^{\mathrm i\alpha_m},
\quad
 v_i(k)=v_s[\Re c_i(k),\Im c_i(k)]^\top,
\end{equation}
\begin{equation}
 q_i(k+1)=\operatorname{mod}(q_i(k)+\Delta t v_i(k),L).
\end{equation}
因此 $v_s$ 是速度尺度而不是恒定速度。单环活动集中度定义为
\begin{equation}
 \kappa_i(k)=\frac{|c_i(k)|}{\sum_m a_{im}(k)+\varepsilon}.
\end{equation}

上述环形连接与活动峰动力学与头方向细胞群的空间朝向表征\upcite{Zhang1996}相一致。自适应环形吸引子模型进一步表明，外部线索可调制以自我为中心的方向编码\upcite{Chen2026}；从连续神经动力学角度，相关研究讨论了预测性神经网络动力学\upcite{Ao2021}，人脑自然导航实验也为周期方向表征提供了神经罗盘证据\upcite{Lu2025}。

原始异中心神经场模型的完整求解流程见算法~\ref{alg:original_model}。

\begin{algorithm}[htbp]
 \caption{原始异中心神经场集群运动}
 \label{alg:original_model}
 \begin{algorithmic}[1]
 \Require 模型参数 $N,M,L,\Delta t,v_s,\beta,\sigma_s,h_t,h_b,\zeta$
 \Require 初始状态 $q_i(0)$、$\phi_i(0)$、$u_i(0)$，或用于随机初始化的随机种子
 \Ensure 轨迹 $\{q_i(k)\}_{i=1,k=0}^{N,K}$、神经状态 $\{u_i(k)\}_{i=1,k=0}^{N,K}$ 和社会输入历史 $\{I_i(k)\}$
 \State 初始化
 $q_i(0)\sim\mathcal U([0,L)^2)$，
 $\phi_i(0)\sim\mathcal U(0,2\pi)$，
 $u_i(0)=0.1\epsilon_i$，
 $\epsilon_i\sim\mathcal N(0,I_M)$
 \State 设置 $\alpha_m=2\pi m/M+\phi_i(0)$，并在后续时刻保持固定
 \For{$k=0,1,\ldots,K-1$}
  \For{所有 $i=1,2,\ldots,N$，并行}
   \State $\mathcal B_i\leftarrow\{j:j\ne i\}$
   \For{$j\in\mathcal B_i$}
    \State 计算周期差分 $r_{ij}(k)$ 和世界方位 $\theta_{ij}(k)$
    \State 将邻居方位累加到圆周高斯输入 $I^{\mathrm{soc}}_{im}(k)$
   \EndFor
   \State $I_i(k)\leftarrow I^{\mathrm{soc}}_{i}(k)$
   \State $u_i(k+1)\leftarrow
   u_i(k)+\Delta t\left[-u_i(k)+\frac{1}{M}W\tanh(\beta u_i(k))-h_b\mathbf 1+I_i(k)\right]$
   \State 计算正活动 $a_{im}(k)$ 和复数合矢量 $c_i(k)$
   \State 解码速度 $v_i(k)=v_s[\Re c_i(k),\Im c_i(k)]^\top$
  \EndFor
  \For{所有 $i=1,2,\ldots,N$，并行}
   \State 更新位置 $q_i(k+1)=\operatorname{mod}(q_i(k)+\Delta t v_i(k),L)$
  \EndFor
 \EndFor
 \State \Return $\{q_i(k)\}$、$\{u_i(k)\}$、$\{I_i(k)\}$
 \end{algorithmic}
\end{algorithm}

算法~\ref{alg:original_model} 对应原始模型的全邻居版本，其中所有个体均使用同一时刻的位置计算输入，并在该时间步结束时同步更新。GO、LO、凝聚距离和活动集中度等指标不参与状态更新，而是由返回的轨迹和神经状态离线计算。

需要注意的是，上述模型采用全邻居、等权交互，而真实群体中的邻居选择可能受视觉线索\upcite{Krongauz2024}、注意力机制\upcite{Li2024}和选择性交互\upcite{Xue2025}影响。本文暂不建模这些机制，而以固定全邻居输入隔离间歇观测对 flocking 的影响。

\subsection{原始模型在间歇观测下需要解决的问题}

原始模型的求解依赖一个强假设：$\mathcal B_i(k)$ 中每个邻居的真实方位 $\theta_{ij}(k)$ 在每一步都可获得。当遮挡或通信中断发生时，该假设不再成立，主要体现在以下几个方面。

\paragraph{社会输入缺失}
若直接删除缺失邻居，则
\begin{equation}
 I_i^{\mathrm{drop}}(k)=
 h_0\sum_{j\in\mathcal B_i(k)}\delta_{ij}(k)
 g_\sigma(\alpha_m-\theta_{ij}(k)).
\end{equation}
输入总量和方位覆盖同时下降。缺失邻居越多，神经场活动峰越弱，甚至可能被剩余邻居方向主导。

\paragraph{过期方位造成错误峰}
若保持最后一次观测方位，则输入幅值可以维持，但方向可能已经过期：
\begin{equation}
 \hat\theta_{ij}(k)=z_{ij}(k_{\mathrm{last}}),
\quad k-k_{\mathrm{last}}>0.
\end{equation}
在群体转向时，过期方位会把神经活动峰持续固定在旧方向，形成“幽灵邻居输入”，使个体无法及时跟随方向变化。

\paragraph{只有方位观测，不能直接估计位置和速度}
传感器不提供距离，因此不能使用二维位置—速度 Kalman 滤波器直接恢复邻居状态。可观测的局部模型只能建立在方位及其时间变化率上。连续不同时刻的方位观测可以识别局部角速度，但单次方位观测不能直接识别角速度。

\paragraph{轨迹生命周期和数据关联问题}
遮挡期间必须区分“暂时没有观测但目标仍可能存在”和“目标已经离开或轨迹失效”。如果轨迹无限保留，会累积错误预测；如果过早删除，重新出现时又需要重新初始化。因此需要观测年龄、预测方差和有效集合 $m_{ij}(k)$ 的管理机制。分布式多传感器多目标滤波为间歇观测下的轨迹维持提供了参考\upcite{Qi2024}；对于纯方位目标，交互式多模型与无迹信息滤波可用于处理轨迹存在性和跟踪不确定性\upcite{Pan2026}。

\paragraph{评价指标可能掩盖退化}
当群体已经进入直线运动时，即使删除部分邻居输入，群体也可能依靠既有航向惯性维持较高 GO。因此，仅报告直线运动下的 GO 不足以评估预测补偿。需要构造遮挡前受控转向实验，并报告相对无遮挡控制的 GO、LO、凝聚距离和恢复误差。

\paragraph{安全约束不在原始模型中}
原始异中心神经场模型没有显式分离、排斥或碰撞屏障项。即使 Kalman 改善了有序度和凝聚恢复，也不能自动获得安全保证。因此本文将最小距离和危险接近比例作为独立的风险诊断指标，而不将有序度改善解释为碰撞安全改善。

\section{间歇方位观测下的Kalman预测与神经场输入补偿}

邻居方位观测的间歇缺失会减少神经场的社会输入，并改变其方向结构。为在有限缺失窗口内延续已观测邻居的方向信息，本文建立由方位状态估计、缺失期间预测和神经场输入重建组成的补偿方法。每个智能体独立维护其已跟踪邻居的方位--方位角速度状态：观测到达时执行 Kalman 更新，观测缺失时保留模型预测，并将预测均值及协方差映射为神经场的周期方向输入。

\subsection{方法概述与观测模型}
设 $\delta_{ij}(k)\in\{0,1\}$ 表示观测是否新鲜到达。若传感器在机体坐标系中输出方位，且自身世界朝向 $\phi_i(k)$ 已知，则
\begin{equation}
 z^b_{ij}(k)=\operatorname{wrap}(\theta_{ij}(k)-\phi_i(k)+\nu_{ij}(k)),
\end{equation}
\begin{equation}
 z_{ij}(k)=\operatorname{wrap}(z^b_{ij}(k)+\phi_i(k)),
 \quad \nu_{ij}(k)\sim\mathcal N(0,R).
\end{equation}
仅当 $\delta_{ij}(k)=1$ 时执行测量更新。本文假设邻居目标编号正确，尚不处理数据关联错误和自身朝向误差。

\subsection{方位状态预测与测量更新}
\subsubsection{方位--方位角速度状态与阻尼扩展}
由于没有距离观测，定义局部展开角状态
\begin{equation}
 x_{ij}(k)=\begin{bmatrix}\vartheta_{ij}(k)\\ \omega_{ij}(k)\end{bmatrix},
 \quad \theta_{ij}(k)=\operatorname{wrap}(\vartheta_{ij}(k)),
\end{equation}
其中 $\vartheta_{ij}(k)$ 是局部展开角，$\theta_{ij}(k)$ 是其对应的圆周方位，$\omega_{ij}(k)$ 是邻居相对于观测者的方位角速度，并非邻居自身航向的角速度。基础模型采用局部常角速度假设
\begin{equation}
 x_{ij}(k+1)=F x_{ij}(k)+w_{ij}(k),
 \quad F=\begin{bmatrix}1&\Delta t\\0&1\end{bmatrix},
\end{equation}
\begin{equation}
 Q=q_\omega\begin{bmatrix}(\Delta t)^3/3&(\Delta t)^2/2\\(\Delta t)^2/2&\Delta t\end{bmatrix},
 \quad z_{ij}(k)=Hx_{ij}(k)+\nu_{ij}(k),\quad H=\begin{bmatrix}1&0\end{bmatrix}.
\end{equation}
为提高长时间缺失下的外推稳定性，进一步采用阻尼角速度模型作为扩展。令
$\rho=\exp(-\Delta t/\tau)$，其中 $\tau$ 为角速度相关时间，则
\begin{equation}
 x_{ij}(k+1)=F_{\mathrm{GM}}x_{ij}(k)+w_{ij}(k),
 \quad
 F_{\mathrm{GM}}=\begin{bmatrix}1&\tau(1-\rho)\\0&\rho\end{bmatrix},
\end{equation}
\begin{equation}
 Q_{\mathrm{GM}}=q_\omega
 \begin{bmatrix}
 \tau^2\Delta t-\tau^3\left(2(1-\rho)-\frac{1-\rho^2}{2}\right)
 & \frac{\tau^2(1-\rho)^2}{2}\\[2mm]
 \frac{\tau^2(1-\rho)^2}{2}
 & \frac{\tau(1-\rho^2)}{2}
 \end{bmatrix}.
\end{equation}
当 $\tau\to\infty$ 时，该模型退化为常角速度模型；当 $\tau$ 较小时，角速度预测向零衰减，从而降低长时间外推对过期角速度的依赖。上述模型只预测方位，不恢复邻居的二维位置或速度。当过程模型参数或噪声比例难以准确给定、且模型失配主要来自参数比例时，可采用自适应 Kalman 滤波思想在线调整噪声强度\upcite{Ge2025}。

\subsubsection{初始化、预测与更新}
首次收到目标观测时初始化
\begin{equation}
 \hat x_{ij}(k)=\begin{bmatrix}z_{ij}(k)\\0\end{bmatrix},
 \quad P_{ij}(k)=\operatorname{diag}(R,P_{\omega,0}).
\end{equation}
对已初始化轨迹先执行
\begin{equation}
 \hat x^-_{ij}(k)=F\hat x_{ij}(k-1),\qquad
 P^-_{ij}(k)=FP_{ij}(k-1)F^\top+Q.
\end{equation}
当 $\delta_{ij}(k)=1$ 时，计算圆周创新
\begin{equation}
 e_{ij}(k)=\operatorname{wrap}\left(z_{ij}(k)-H\hat x^-_{ij}(k)\right),
\end{equation}
\begin{equation}
 S_{ij}=HP^-_{ij}H^\top+R,\quad K_{ij}=P^-_{ij}H^\top S_{ij}^{-1},
\end{equation}
\begin{equation}
 \hat x_{ij}(k)=\hat x^-_{ij}(k)+K_{ij}e_{ij}(k),
\end{equation}
\begin{equation}
 P_{ij}(k)=(I-K_{ij}H)P^-_{ij}(k)(I-K_{ij}H)^\top+K_{ij}RK_{ij}^\top.
\end{equation}
对于少量异常方位观测，还可借鉴最大似然最大相关熵准则抑制非高斯重尾噪声对状态更新的影响\upcite{ZhangBearing2025}；本文主实验仍采用高斯测量模型，以便与原始神经场输入保持一致。

当 $\delta_{ij}(k)=0$ 时跳过更新，仅保留预测值。记观测年龄为 $a_{ij}(k)$，成功更新时置零，连续缺失时每步增加 $\Delta t$。当 $a_{ij}(k)>\tau_{\max}$ 时删除轨迹，后续重新获得观测时重新初始化。记轨迹有效标记为 $m_{ij}(k)\in\{0,1\}$，其中观测指示量 $\delta_{ij}(k)$ 控制是否执行测量更新，$m_{ij}(k)$ 控制是否向神经场注入该邻居的方向输入。这种在随机丢包期间保留预测、恢复后重新利用轨迹的处理与网络化系统中的丢包补偿思想一致\upcite{Yu2024}。在分布式实现中，还可采用自适应通信策略降低信息交换需求\upcite{Selvi2025}，或在随机通信协议下通过一致性滤波保持节点估计一致\upcite{Zhou2025}；融合多源测量协方差能够进一步提高共同估计精度\upcite{Yang2025}。

观测矩阵为
\begin{equation}
 \mathcal O=\begin{bmatrix}H\\HF\end{bmatrix}=\begin{bmatrix}1&0\\1&\Delta t\end{bmatrix}.
\end{equation}
当 $\Delta t\ne0$ 时 $\det(\mathcal O)=\Delta t\ne0$，说明连续不同时间的方位观测可识别局部二状态模型；单次方位观测不能直接识别角速度。若方位观测与状态之间的关系出现更强非线性，可进一步采用无迹 Kalman 滤波类方法处理均值与协方差的非线性传播\upcite{Rong2024}。

\subsection{不确定性感知神经场输入}

\subsubsection{预测均值输入}
定义
\begin{equation}
 \mu_{ij}(k)=H\hat x_{ij}(k),\quad V_{ij}(k)=HP_{ij}(k)H^\top.
\end{equation}
若轨迹已经初始化且未过期，记 $m_{ij}(k)=1$，否则为零。最小预测均值输入为
\begin{equation}
 I^{\mathrm{mean}}_{im}(k)=h_0\sum_jm_{ij}(k)
 g_\sigma(\alpha_m-\mu_{ij}(k)),
\end{equation}
其中 $g_\sigma$ 为式（4）中的圆周高斯核。需要区分：$\delta_{ij}$ 控制是否更新滤波器，$m_{ij}$ 控制预测轨迹是否继续注入神经场。

\subsubsection{协方差传播与周期卷积}
在局部高斯假设下，方位后验为 wrapped normal 分布
\begin{equation}
 \Theta_{ij}(k)\mid\mathcal Y_{i,k}\sim\mathrm{WN}(\mu_{ij}(k),V_{ij}(k)).
\end{equation}
将社会输入定义为方位核的条件期望：
\begin{equation}
 \bar g_\sigma(x;V)=\int_{-\pi}^{\pi}g_\sigma(x-\epsilon)p_{\mathrm{WN}}(\epsilon;0,V)\,d\epsilon,
\end{equation}
\begin{equation}
 I^{\mathrm{unc}}_{im}(k)=h_0\sum_jm_{ij}(k)\bar g_\sigma(\alpha_m-\mu_{ij}(k);V_{ij}(k)).
\end{equation}

令
\begin{equation}
 c_n=\frac{1}{2\pi}\int_{-\pi}^{\pi}g_\sigma(x)e^{-\mathrm inx}\,dx,
\end{equation}
则周期卷积的 Fourier 形式为
\begin{equation}
 \bar g_\sigma(x;V)=\sum_{n=-\infty}^{\infty}c_n e^{-n^2V/2}e^{\mathrm inx}.
\end{equation}
因此第 $n$ 阶方向信息保留比例为 $e^{-n^2V/2}$，一阶方向信息保留比例为
\begin{equation}
 \gamma(V)=e^{-V/2}.
\end{equation}
当 $V$ 增大时，输入方向结构变宽并趋向均匀背景；由于零阶 Fourier 分量保持不变，不能简单地说不确定性使总输入消失。

在方位分布足够集中时，可采用便于实现的局部高斯近似
\begin{equation}
 \bar g_\sigma(x;V)\approx
 \sqrt{\frac{\sigma_s^2}{\sigma_s^2+V}}
 \exp\left[-\frac{\operatorname{wrap}(x)^2}{2(\sigma_s^2+V)}\right].
\end{equation}
该式只适用于小方差和单一展开分支；大方差情形应使用周期积分或 Fourier 形式。

\subsection{算法流程与伪代码}

每个时间步首先接收并转换方位观测，然后对各邻居轨迹进行初始化、预测或测量更新。经过有效性检查后，轨迹的预测均值和协方差被映射为社会输入，最后驱动原神经场与运动模型。各智能体使用同一时刻的状态同步推进，避免顺序更新产生额外的信息差异。

\begin{algorithm}[htbp]
 \caption{间歇方位观测下的Kalman预测与神经场输入补偿}
 \label{alg:kalman_compensation}
 \scriptsize
 \begin{algorithmic}[1]
 \Require 时间步长 $\Delta t$、测量噪声方差 $R$、初始角速度方差 $P_{\omega,0}$
 \Require 状态模型 $\mathrm{model}\in\{\mathrm{CV},\mathrm{GM}\}$、输入模式 $\mathrm{mode}\in\{\mathrm{mean},\mathrm{uncertainty}\}$
 \Require 过程噪声参数 $q_\omega$、相关时间 $\tau$、轨迹生命周期 $\tau_{\max}$
 \Require 方向单元 $\alpha_m$、核宽度 $\sigma_s$、单边输入幅值 $h_0$ 及原神经场参数
 \Ensure 轨迹集合 $\mathcal T_i$、社会输入 $I_i(k)$ 和同步更新的神经场与运动状态
 \State 初始化所有轨迹集合 $\mathcal T_i\leftarrow\varnothing$
 \State 根据 $\mathrm{model}$ 构造状态矩阵 $F$ 与过程噪声协方差 $Q$
 \If{$\mathrm{mode}=\mathrm{uncertainty}$}
  \State 预计算方向核 Fourier 系数 $c_0,\ldots,c_{N_F}$
 \EndIf
 \end{algorithmic}
\end{algorithm}
\clearpage

\begin{algorithm}[htbp]
 \ContinuedFloat
 \caption{（续）：输入构造及闭环推进}
 \scriptsize
 \begin{algorithmic}
 \For{每个时间步 $k$}
  \State 从同一时刻的系统状态获取所有新鲜方位观测
  \For{每个智能体 $i$}
   \State $I_i(k)\leftarrow\mathbf 0$；$\mathcal J_i\leftarrow$ 已维护轨迹与当前观测邻居身份之并集
   \For{每个 $j\in\mathcal J_i$}
    \If{$\delta_{ij}(k)=1$}
     \State 将机体系观测转换为世界方位 $z_{ij}(k)$
    \EndIf
    \If{$j\notin\mathcal T_i$}
     \If{$\delta_{ij}(k)=0$}
      \State 跳过该邻居
     \Else
      \State $\hat x_{ij}\leftarrow[z_{ij}(k),0]^\top$；$P_{ij}\leftarrow\operatorname{diag}(R,P_{\omega,0})$；$a_{ij}\leftarrow0$
      \State 将 $j$ 加入 $\mathcal T_i$
     \EndIf
    \Else
     \State $\hat x^-_{ij}\leftarrow F\hat x_{ij}$；$P^-_{ij}\leftarrow FP_{ij}F^\top+Q$
     \If{$\delta_{ij}(k)=1$}
      \State $e_{ij}\leftarrow\operatorname{wrap}(z_{ij}(k)-H\hat x^-_{ij})$
      \State $S_{ij}\leftarrow HP^-_{ij}H^\top+R$；$K_{ij}\leftarrow P^-_{ij}H^\top S_{ij}^{-1}$
      \State $\hat x_{ij}\leftarrow\hat x^-_{ij}+K_{ij}e_{ij}$；$a_{ij}\leftarrow0$
      \State $P_{ij}\leftarrow(I_2-K_{ij}H)P^-_{ij}(I_2-K_{ij}H)^\top+K_{ij}RK_{ij}^\top$
     \Else
      \State $\hat x_{ij}\leftarrow\hat x^-_{ij}$；$P_{ij}\leftarrow P^-_{ij}$；$a_{ij}\leftarrow a_{ij}+\Delta t$
     \EndIf
     \If{$a_{ij}>\tau_{\max}$}
      \State 从 $\mathcal T_i$ 删除 $j$ 并跳过该邻居
     \EndIf
     \State $\mu_{ij}\leftarrow H\hat x_{ij}$；$V_{ij}\leftarrow HP_{ij}H^\top$
     \For{每个方向单元 $m$}
      \State $x\leftarrow\operatorname{wrap}(\alpha_m-\mu_{ij})$
      \If{$\mathrm{mode}=\mathrm{mean}$}
       \State $b\leftarrow\exp[-x^2/(2\sigma_s^2)]$
      \Else
       \State $b\leftarrow c_0+\sum_{n=1}^{N_F}2c_n\exp(-n^2V_{ij}/2)\cos(nx)$
      \EndIf
      \State $I_{im}(k)\leftarrow I_{im}(k)+h_0b$
     \EndFor
    \EndIf
   \EndFor
   \State 将 $I_i(k)$ 代入原模型总输入并按相同时间索引推进神经场、解码速度和更新位置
  \EndFor
  \State 进入下一个时间步
 \EndFor
 \State \Return $\{q_i(k)\}$、$\{u_i(k)\}$、$\{I_i(k)\}$
 \end{algorithmic}
\end{algorithm}
\clearpage

通过状态模型与输入模式的组合，可得到三种预测补偿方法：KF-mean 使用常角速度模型的预测均值方向核；KF-uncertainty 使用常角速度模型的协方差卷积方向核；KF-GM-uncertainty 使用阻尼角速度模型和协方差卷积方向核。fresh-drop 与 hold-last 作为实验基线，分别使用当前新鲜方位与最后一次观测方位，不需要执行上述 Kalman 递推。

\section{预测不确定性与神经场输入的理论性质}

间歇方位观测下的异中心神经场需要同时处理方向估计、预测不确定性和社会输入缺失。为刻画这些因素之间的关系，本章以方位后验、周期社会输入和神经场状态为理论对象，建立一条可检验的误差传播链：
\[
\text{方位后验}\;\longrightarrow\;\text{社会输入误差}\;\longrightarrow\;\text{神经场状态误差}\;\longrightarrow\;\text{速度与位置偏差}.
\]

理论分析包含四个相互衔接的层次：首先证明协方差感知输入在相同有效轨迹集合上的条件均方最优性，并量化只使用预测均值所产生的额外重建风险；其次区分完整邻居集合与有效轨迹集合，分解方向预测误差、漏邻居误差和错误关联误差；随后将连续缺失期间的方位方差传播转化为周期方向谐波衰减和方向信息保持窗口；最后在神经场推进映射满足局部 Lipschitz 条件时，给出社会输入扰动到神经场状态、速度解码和位置偏差的有限时间传播界。上述结论分别描述输入重建、方向信息保持和有限时间状态偏差，不将局部输入最优性等同于全局闭环稳定性。

\subsection{理论对象、符号与适用条件}

考虑智能体 $i$ 在时刻 $k$ 的观测历史 $\mathcal Y_i(k)$。记原始模型中的有向邻居集合为 $\mathcal B_i(k)$，由历史观测确定、已经初始化且尚未过期的有效轨迹集合为
\[
\mathcal S_i(k)\subseteq\mathcal B_i(k).
\]
前者表示参考社会输入可能包含的全部邻居，后者表示当前预测器实际可以继续维护的轨迹。理论分析必须区分这两个集合：协方差感知输入只能重建 $\mathcal S_i(k)$ 内的邻居，不能自动恢复从未初始化或已经删除的邻居。

对环形神经场的方向单元 $\alpha_m$，定义周期方位核
\begin{equation}
 g_\sigma(x)=\exp\!\left[-\frac{\operatorname{wrap}(x)^2}{2\sigma_s^2}\right],
 \qquad
 \operatorname{wrap}(x)=\operatorname{mod}(x+\pi,2\pi)-\pi .
 \label{eq:v2_kernel}
\end{equation}
其中 $0<\sigma_s\leq\pi$。圆周距离记为
\begin{equation}
 d_{\Sone}(a,b)=|\operatorname{wrap}(a-b)|.
 \label{eq:v2_circle_distance}
\end{equation}

对 $j\in\mathcal B_i(k)$，令 $\Theta_{ij}(k)\in\Sone$ 表示真实世界方位；对 $j\in\mathcal S_i(k)$，方位预测器给出预测均值 $\mu_{ij}(k)$ 和方位方差 $V_{ij}(k)$。在当前观测历史条件下，假定有效轨迹的方位边缘后验满足
\begin{equation}
 \Theta_{ij}\mid\mathcal Y_i\sim\operatorname{WN}(\mu_{ij},V_{ij}),
 \qquad j\in\mathcal S_i,
 \label{eq:v2_wn_assumption}
\end{equation}
其中 $\operatorname{WN}$ 表示包裹正态分布。不同邻居之间不要求条件独立，因为本章的条件期望和误差分解只使用线性性与三角不等式。该假设还要求局部展开角分支在当前时间窗口内有效、邻居身份已经正确关联，且轨迹仍满足生命周期和有效性规则。

定义有效轨迹集合上的理想社会输入
\begin{equation}
 J_{im}^{(\mathcal S)}=h_0\sum_{j\in\mathcal S_i}g_\sigma(\alpha_m-\Theta_{ij}),
 \qquad
 J_i^{(\mathcal S)}=[J_{i0}^{(\mathcal S)},\ldots,J_{i,M-1}^{(\mathcal S)}]^{\top}.
 \label{eq:v2_tracked_ideal}
\end{equation}
完整邻居集合上的新鲜参考输入定义为
\begin{equation}
 J_{im}^{(\mathcal B)}=h_0\sum_{j\in\mathcal B_i}g_\sigma(\alpha_m-\Theta_{ij}),
 \qquad
 J_i^{(\mathcal B)}=[J_{i0}^{(\mathcal B)},\ldots,J_{i,M-1}^{(\mathcal B)}]^{\top}.
 \label{eq:v2_full_ideal}
\end{equation}
显然，$J_i^{(\mathcal B)}$ 与 $J_i^{(\mathcal S)}$ 只有在当前有效轨迹集合覆盖全部参考邻居时才一致。

令
\begin{equation}
 \bar g_\sigma(x;V)=\int_{-\pi}^{\pi}g_\sigma(x-\epsilon)
 p_{\operatorname{WN}}(\epsilon;0,V)\,\mathrm d\epsilon,
 \label{eq:v2_conv_kernel}
\end{equation}
并定义协方差感知输入
\begin{equation}
 \bar J_{im}^{(\mathcal S)}=h_0\sum_{j\in\mathcal S_i}
 \bar g_\sigma(\alpha_m-\mu_{ij};V_{ij}),
 \qquad
 \bar J_i^{(\mathcal S)}=[\bar J_{i0}^{(\mathcal S)},\ldots,\bar J_{i,M-1}^{(\mathcal S)}]^{\top}.
 \label{eq:v2_unc_input}
\end{equation}

\subsection{协方差感知输入的条件均方最优性}

\begin{theorem}[有效轨迹集合上的条件均方最优性]
\label{thm:v2_mmse}
在式\eqref{eq:v2_wn_assumption}的后验模型、相同有效轨迹集合 $\mathcal S_i$ 和给定观测历史 $\mathcal Y_i$ 下，式\eqref{eq:v2_unc_input}满足
\begin{equation}
 \bar J_i^{(\mathcal S)}=\E[J_i^{(\mathcal S)}\mid\mathcal Y_i].
 \label{eq:v2_conditional_mean}
\end{equation}
因此，对任意仅依赖 $\mathcal Y_i$ 且具有有限二阶矩的输入估计 $\widetilde J_i$，有
\begin{equation}
 \begin{aligned}
 \E[\|J_i^{(\mathcal S)}-\widetilde J_i\|_2^2\mid\mathcal Y_i]
 ={}&\E[\|J_i^{(\mathcal S)}-\bar J_i^{(\mathcal S)}\|_2^2\mid\mathcal Y_i]\\
 &+\|\bar J_i^{(\mathcal S)}-\widetilde J_i\|_2^2.
 \end{aligned}
 \label{eq:v2_orthogonal_projection}
\end{equation}
故协方差感知输入在该估计类中唯一地最小化有效轨迹集合上的理想社会输入条件均方重建误差，除去概率为零的等价情形。
\end{theorem}

\begin{proof}
对任意方向单元 $\alpha_m$，由条件期望线性性和式\eqref{eq:v2_wn_assumption}，有
\begin{equation}
 \begin{aligned}
 \E[J_{im}^{(\mathcal S)}\mid\mathcal Y_i]
 &=h_0\sum_{j\in\mathcal S_i}
 \E[g_\sigma(\alpha_m-\Theta_{ij})\mid\mathcal Y_i]\\
 &=h_0\sum_{j\in\mathcal S_i}
 \bar g_\sigma(\alpha_m-\mu_{ij};V_{ij})
 =\bar J_{im}^{(\mathcal S)}.
 \end{aligned}
 \label{eq:v2_component_expectation}
\end{equation}
逐分量成立即得式\eqref{eq:v2_conditional_mean}。再将误差分解为
\begin{equation}
 J_i^{(\mathcal S)}-\widetilde J_i
 =\bigl(J_i^{(\mathcal S)}-\bar J_i^{(\mathcal S)}\bigr)
 +\bigl(\bar J_i^{(\mathcal S)}-\widetilde J_i\bigr).
 \label{eq:v2_error_decomposition}
\end{equation}
第二项由 $\mathcal Y_i$ 确定，且第一项的条件均值为零，因此平方展开后的交叉项条件期望为零，得到式\eqref{eq:v2_orthogonal_projection}。
\end{proof}

定理\ref{thm:v2_mmse}不是 Kalman 滤波器递推本身的新最优性，而是一个针对异中心神经场输入接口的投影结论：当预测器提供方位后验时，把后验分布作用于方向核所得的输入，正是有效轨迹集合上理想社会输入的条件期望。因此，协方差不再只是滤波器内部的诊断量，而成为神经场输入构造的一部分。

\subsection{均值输入的额外重建风险}

KF-mean 只把预测均值代入方向核：
\begin{equation}
 J_{im}^{\mathrm{mean}}=h_0\sum_{j\in\mathcal S_i}
 g_\sigma(\alpha_m-\mu_{ij}).
 \label{eq:v2_mean_input}
\end{equation}

\begin{corollary}[忽略方位协方差的额外风险]
\label{cor:v2_mean_risk}
在定理\ref{thm:v2_mmse}的条件下，KF-mean 相对于协方差感知输入增加的条件均方重建风险为
\begin{equation}
 \E[\|J_i^{(\mathcal S)}-J_i^{\mathrm{mean}}\|_2^2\mid\mathcal Y_i]
 -\E[\|J_i^{(\mathcal S)}-\bar J_i^{(\mathcal S)}\|_2^2\mid\mathcal Y_i]
 =\|\bar J_i^{(\mathcal S)}-J_i^{\mathrm{mean}}\|_2^2\ge0.
 \label{eq:v2_mean_risk_gap}
\end{equation}
当所有 $V_{ij}\to0$ 时，$\bar J_i^{(\mathcal S)}$ 与 $J_i^{\mathrm{mean}}$ 趋于一致；在非零方差下，二者一般不一致。
\end{corollary}

\begin{proof}
在式\eqref{eq:v2_orthogonal_projection}中取 $\widetilde J_i=J_i^{\mathrm{mean}}$ 即得。
\end{proof}

对 $0<\sigma_s\le\pi$，周期核关于圆周距离满足 Lipschitz 界
\begin{equation}
 |g_\sigma(a)-g_\sigma(b)|
 \le L_\sigma d_{\Sone}(a,b),
 \qquad L_\sigma=\frac{e^{-1/2}}{\sigma_s}.
 \label{eq:v2_lipschitz_kernel}
\end{equation}
在局部展开角分支有效且 $\E[d_{\Sone}(\mu_{ij},\Theta_{ij})^2\mid\mathcal Y_i]\le V_{ij}$ 时，有
\begin{equation}
 \E[\|J_i^{\mathrm{mean}}-J_i^{(\mathcal S)}\|_\infty\mid\mathcal Y_i]
 \le h_0L_\sigma\sum_{j\in\mathcal S_i}\sqrt{V_{ij}}.
 \label{eq:v2_mean_bound}
\end{equation}
该界不包含漏掉邻居或错误关联造成的误差，因而只描述已跟踪集合内的均值预测风险。

\subsection{完整社会输入的误差分解}

\begin{proposition}[预测误差与漏邻居误差的分解]
\label{prop:v2_full_decomposition}
令 $\mathcal S_i\subseteq\mathcal B_i$，则有
\begin{equation}
 J_i^{(\mathcal B)}-\bar J_i^{(\mathcal S)}
 =\bigl(J_i^{(\mathcal S)}-\bar J_i^{(\mathcal S)}\bigr)
 +J_i^{(\mathcal B\setminus\mathcal S)}.
 \label{eq:v2_full_decomposition}
\end{equation}
进一步，在 $0\le g_\sigma\le1$ 时，满足
\begin{equation}
 \|J_i^{(\mathcal B)}-\bar J_i^{(\mathcal S)}\|_\infty
 \le
 \|J_i^{(\mathcal S)}-\bar J_i^{(\mathcal S)}\|_\infty
 +h_0|\mathcal B_i\setminus\mathcal S_i|.
 \label{eq:v2_full_bound}
\end{equation}
若再满足式\eqref{eq:v2_mean_bound}的条件，则由条件 Jensen 不等式和三角不等式，
\begin{equation}
 \begin{aligned}
 &\E[\|J_i^{(\mathcal S)}-\bar J_i^{(\mathcal S)}\|_\infty\mid\mathcal Y_i]\\
 &\quad\le
 2\E[\|J_i^{(\mathcal S)}-J_i^{\mathrm{mean}}\|_\infty\mid\mathcal Y_i]
 \le2h_0L_\sigma\sum_{j\in\mathcal S_i}\sqrt{V_{ij}}.
 \end{aligned}
 \label{eq:v2_unc_absolute_bound}
\end{equation}
进而有
\begin{equation}
 \E[\|J_i^{(\mathcal B)}-\bar J_i^{(\mathcal S)}\|_\infty\mid\mathcal Y_i]
 \le h_0\left(
 2L_\sigma\sum_{j\in\mathcal S_i}\sqrt{V_{ij}}
 +|\mathcal B_i\setminus\mathcal S_i|
 \right).
 \label{eq:v2_total_input_bound}
\end{equation}
\end{proposition}

\begin{proof}
将式\eqref{eq:v2_tracked_ideal}和式\eqref{eq:v2_full_ideal}按集合 $\mathcal S_i$ 与 $\mathcal B_i\setminus\mathcal S_i$ 拆分即可得到式\eqref{eq:v2_full_decomposition}。由于 $g_\sigma$ 与 $\bar g_\sigma$ 的值均不超过 1，遗漏集合对每一个方向单元的输入幅值不超过 $h_0|\mathcal B_i\setminus\mathcal S_i|$，再使用三角不等式得到式\eqref{eq:v2_full_bound}。另一方面，$\bar J_i^{(\mathcal S)}=\E[J_i^{(\mathcal S)}\mid\mathcal Y_i]$，故
\[
 \|J_i^{\mathrm{mean}}-\bar J_i^{(\mathcal S)}\|_\infty
 \le\E[\|J_i^{\mathrm{mean}}-J_i^{(\mathcal S)}\|_\infty\mid\mathcal Y_i].
\]
将其与三角不等式及式\eqref{eq:v2_mean_bound}结合，得到式\eqref{eq:v2_unc_absolute_bound}；最后代入式\eqref{eq:v2_full_bound}得到式\eqref{eq:v2_total_input_bound}。
\end{proof}

命题\ref{prop:v2_full_decomposition}区分了两种不同失效机制：第一项是已经建立轨迹但预测不准，第二项是邻居没有进入有效轨迹集合。协方差卷积只能降低第一项中的方向重建风险，不能凭空恢复第二项。因此，协方差感知输入的条件均方最优性不能被解释为完整邻居集合上的无条件最优性。

若需要描述错误数据关联，可另定义错误轨迹集合 $\mathcal A_i$，并将其输入作为第三项加入式\eqref{eq:v2_full_decomposition}。本章不对该项建立概率模型，只保留其作为模型适用边界。

\subsection{方向谐波衰减与信息保持窗口}

令原方向核的 Fourier 系数为
\begin{equation}
 c_n=\frac{1}{2\pi}\int_{-\pi}^{\pi}g_\sigma(x)e^{-\mathrm inx}\,\mathrm dx .
 \label{eq:v2_fourier_coeff}
\end{equation}
由于 wrapped normal 的第 $n$ 阶特征函数为 $\exp(-n^2V/2)$，周期卷积核的系数为
\begin{equation}
 c_n(V)=c_n\exp\left(-\frac{n^2V}{2}\right).
 \label{eq:v2_harmonic_decay}
\end{equation}
因此一阶方向信息的保留比例为
\begin{equation}
 \gamma_1(V)=\exp(-V/2).
 \label{eq:v2_first_harmonic}
\end{equation}

在常角速度模型下，设最后一次观测更新后的协方差为
\begin{equation}
 P_0=\begin{bmatrix}p_{11}&p_{12}\\p_{12}&p_{22}\end{bmatrix}.
 \label{eq:v2_p0}
\end{equation}
连续缺失 $T=b\Delta t$ 后，方位预测方差为
\begin{equation}
 V(T)=p_{11}+2Tp_{12}+T^2p_{22}+\frac{q_\omega T^3}{3}.
 \label{eq:v2_vt}
\end{equation}

\begin{corollary}[一阶方向信息保持窗口]
\label{cor:v2_direction_window}
给定一阶方向信息的最低保留比例 $\gamma_{\min}\in(0,1)$，令
\begin{equation}
 V_{\max}=-2\ln\gamma_{\min}.
 \label{eq:v2_vmax}
\end{equation}
则理论方向信息窗口定义为
\begin{equation}
 T_\gamma=\sup\left\{T\ge0:
 V(s)\le V_{\max},\ \forall s\in[0,T]\right\}.
 \label{eq:v2_tgamma}
\end{equation}
在 $[0,T_\gamma]$ 内有 $\gamma_1(s)\ge\gamma_{\min}$。当 $p_{12}<0$ 时，$V(s)$ 在极短时间内可能不单调，因此必须检查整个区间而不能只检查终点。
\end{corollary}

\begin{proof}
由式\eqref{eq:v2_first_harmonic}，$\gamma_1(s)\ge\gamma_{\min}$ 当且仅当 $V(s)\le-2\ln\gamma_{\min}=V_{\max}$。对所有 $s\in[0,T]$同时要求该条件，即得到式\eqref{eq:v2_tgamma}。
\end{proof}

若只需要单调的保守上界，可使用
\begin{equation}
 \overline V(T)=p_{11}+2T|p_{12}|+T^2p_{22}+\frac{q_\omega T^3}{3},
 \qquad V(T)\le\overline V(T),
 \label{eq:v2_vupper}
\end{equation}
并以 $\overline V(T)=V_{\max}$ 的最小正根定义保守窗口。该窗口是单条已跟踪邻居的方向信息窗口，不是集群稳定时间、碰撞安全时间或邻居存在时间。

\subsection{神经场的有限时间输入—状态误差界}

为刻画社会输入误差对神经场动力学的影响，定义神经场的无输入推进映射
\begin{equation}
 f(u)=-u+\frac{1}{M}W\tanh(\beta u)-h_b\bm 1,
 \qquad
 \Phi(u)=u+\Delta t f(u).
 \label{eq:v2_field_map}
\end{equation}
由于逐分量的 $\tanh(\beta u)$ 是 $\beta$-Lipschitz，推进映射在全空间上可取保守常数
\begin{equation}
 L_\Phi^{\mathrm{glob}}
 =|1-\Delta t|+\frac{\Delta t\beta}{M}\|W\|_2.
 \label{eq:v2_global_lipschitz}
\end{equation}
若已知两条轨迹均落在状态区域 $\mathcal U$，则还可以通过 Jacobian 得到更紧的局部常数
\begin{equation}
 L_\Phi(\mathcal U)=
 \sup_{u\in\mathcal U}
 \left\|(1-\Delta t)I_M+
 \frac{\Delta t\beta}{M}W
 \operatorname{diag}\!\left(\operatorname{sech}^2(\beta u)\right)
 \right\|_2.
 \label{eq:v2_local_lipschitz}
\end{equation}
全局常数适合给出无需轨迹信息的保守界；局部常数可利用神经元已经进入饱和区时 $\operatorname{sech}^2(\beta u)$ 较小的事实，但必须由给定状态区域或保存轨迹独立验证。
考虑从同一检查点开始的两条轨迹：fresh 轨迹使用 $J_i^{(\mathcal B)}$，预测轨迹使用 $\widehat J_i$。令
\begin{equation}
 e_{u,i}(k)=u_i^{\mathrm{pred}}(k)-u_i^{\mathrm{fresh}}(k),
 \qquad
 \Delta I_i(k)=\widehat J_i(k)-J_i^{(\mathcal B)}(k).
 \label{eq:v2_state_error}
\end{equation}

\begin{theorem}[神经场状态的有限时间输入—状态误差界]
\label{thm:v2_field_bound}
设两条轨迹在区间 $k\in\{k_0,\ldots,k_0+H\}$ 内均落在状态区域 $\mathcal U$，且神经场推进映射满足
\begin{equation}
 \|\Phi(u)-\Phi(v)\|_2\le L_\Phi\|u-v\|_2,
 \qquad u,v\in\mathcal U.
 \label{eq:v2_lipschitz_phi}
\end{equation}
则对任意 $\ell\in\{0,\ldots,H\}$，有
\begin{equation}
 \|e_{u,i}(k_0+\ell)\|_2
 \le L_\Phi^\ell\|e_{u,i}(k_0)\|_2
 +\Delta t\sum_{r=0}^{\ell-1}
 L_\Phi^{\ell-1-r}\|\Delta I_i(k_0+r)\|_2.
 \label{eq:v2_field_finite_bound}
\end{equation}
特别地，当两条轨迹从同一检查点分叉，即 $e_{u,i}(k_0)=0$，状态偏差完全由输入误差传播项决定。若进一步有 $L_\Phi<1$ 且 $\|\Delta I_i(k)\|_2\le\varepsilon_I$，则
\begin{equation}
 \|e_{u,i}(k_0+\ell)\|_2
 \le
 \frac{\Delta t(1-L_\Phi^\ell)}{1-L_\Phi}\varepsilon_I
 \le\frac{\Delta t}{1-L_\Phi}\varepsilon_I.
 \label{eq:v2_field_contracting_bound}
\end{equation}
\end{theorem}

\begin{proof}
由两条神经场递推式相减，得到
\begin{equation}
 e_{u,i}(k+1)=\Phi(u_i^{\mathrm{pred}}(k))-
 \Phi(u_i^{\mathrm{fresh}}(k))+\Delta t\Delta I_i(k).
 \end{equation}
对上式应用式\eqref{eq:v2_lipschitz_phi}和三角不等式，得到
\begin{equation}
 \|e_{u,i}(k+1)\|_2
 \le L_\Phi\|e_{u,i}(k)\|_2+\Delta t\|\Delta I_i(k)\|_2.
\end{equation}
反复递推即得式\eqref{eq:v2_field_finite_bound}。当 $L_\Phi<1$ 时，对等比和求和即可得到式\eqref{eq:v2_field_contracting_bound}。
\end{proof}

条件\eqref{eq:v2_lipschitz_phi}是局部状态区域条件，不等价于当前神经场参数在全空间上的全局收缩性。特别是 $\beta$ 较大时，不能未经验证就令 $L_\Phi<1$。因此论文中的严谨结论应分为两层：无收缩假设时使用有限时间界\eqref{eq:v2_field_finite_bound}；只有在给定状态区域的局部收缩条件被单独验证后，才使用有界式\eqref{eq:v2_field_contracting_bound}。

将命题\ref{prop:v2_full_decomposition}代入定理\ref{thm:v2_field_bound}，可得遮挡期间神经场状态偏差由两部分组成：已跟踪邻居的方向预测误差，以及被遗漏邻居带来的输入缺口。该分解解释了为什么协方差卷积可以减小过期方向输入的重建风险，却不能替代邻居存在性判断或数据关联机制。

\subsection{速度解码与位置偏差的有限时间界}

神经场输出的正活动为
\begin{equation}
 a_{im}=\max\{\tanh(\beta u_{im}),0\},
 \qquad
 c_i=\sum_{m=0}^{M-1}a_{im}e^{\mathrm i\alpha_m},
 \qquad
 v_i=v_s[\Re c_i,\Im c_i]^{\top}.
 \label{eq:v2_velocity_decode}
\end{equation}

\begin{lemma}[神经场状态到速度的 Lipschitz 界]
\label{lem:v2_velocity_bound}
对任意两个神经场状态 $u$ 和 $v$，有
\begin{equation}
 \|v(u)-v(v)\|_2
 \le v_s\beta\sqrt M\|u-v\|_2.
 \label{eq:v2_velocity_bound}
\end{equation}
\end{lemma}

\begin{proof}
函数 $x\mapsto\max\{\tanh(\beta x),0\}$ 是 $\beta$-Lipschitz，因此正活动向量满足
\(\|a(u)-a(v)\|_2\le\beta\|u-v\|_2\)。复数合矢量是 $M$ 个单位模方向向量的线性组合，故
\(\|c(u)-c(v)\|_2\le\sqrt M\|a(u)-a(v)\|_2\)。乘以速度尺度 $v_s$ 即得式\eqref{eq:v2_velocity_bound}。
\end{proof}

令 $d_{\Ttwo}$ 表示二维周期区域上的最小镜像距离。由于周期取模映射在该距离下不扩张，两条轨迹满足
\begin{equation}
 d_{\Ttwo}(q_i^{\mathrm{pred}}(k+1),q_i^{\mathrm{fresh}}(k+1))
 \le d_{\Ttwo}(q_i^{\mathrm{pred}}(k),q_i^{\mathrm{fresh}}(k))
 +\Delta t\|v_i^{\mathrm{pred}}(k)-v_i^{\mathrm{fresh}}(k)\|_2.
 \label{eq:v2_position_step}
\end{equation}
结合引理\ref{lem:v2_velocity_bound}和定理\ref{thm:v2_field_bound}，从同一检查点开始的有限时间位置偏差满足
\begin{equation}
 \begin{aligned}
 d_{\Ttwo}(q_i^{\mathrm{pred}}(k_0+\ell),q_i^{\mathrm{fresh}}(k_0+\ell))
 \le{}&v_s\beta\sqrt M\Delta t\\
 &\times\sum_{r=0}^{\ell-1}\|e_{u,i}(k_0+r)\|_2.
 \end{aligned}
 \label{eq:v2_position_bound}
\end{equation}
该式是个体有限时间位置偏差界，不是集群 GO、LO 或碰撞概率的直接上界。后者还需要对速度归一化、个体间耦合和危险距离阈值增加额外假设，因此本章不将式\eqref{eq:v2_position_bound}扩展为闭环安全定理。

\subsection{三类条件有效窗口}

为避免把单一方向指标误称为集群稳定时间，本章区分三类窗口。

\begin{definition}[方向信息窗口]
\label{def:v2_direction_window}
给定 $\gamma_{\min}$，式\eqref{eq:v2_tgamma}定义的一阶谐波保留窗口称为方向信息窗口 $T_\gamma$。
\end{definition}

\begin{definition}[社会输入误差窗口]
\label{def:v2_input_window}
给定输入容许误差 $\varepsilon_I$，定义
\begin{equation}
 T_I(\varepsilon_I)=\sup\left\{T\ge0:
 \sup_{0\le s\le T}\mathcal E_I(s)\le\varepsilon_I\right\},
 \label{eq:v2_input_window}
\end{equation}
其中 $\mathcal E_I(s)$ 可以取式\eqref{eq:v2_total_input_bound}右侧的条件期望上界，或由校准后的预测误差模型给出。
\end{definition}

\begin{definition}[神经场状态窗口]
\label{def:v2_field_window}
给定神经场状态容许误差 $\varepsilon_u$，定义
\begin{equation}
 T_u(\varepsilon_u)=\sup\left\{T\ge0:
 \sup_{0\le s\le T}\mathcal E_u(s)\le\varepsilon_u\right\},
 \label{eq:v2_field_window}
\end{equation}
其中 $\mathcal E_u(s)$ 由式\eqref{eq:v2_field_finite_bound}计算。
\end{definition}

实际预测器的可用窗口还受到轨迹年龄截止和方位方差截止约束。因此可定义条件可用窗口
\begin{equation}
 T_{\mathrm{usable}}
 =\min\left\{T_\gamma,
 T_I(\varepsilon_I),
 T_u(\varepsilon_u),
 \tau_{\max}\right\}.
 \label{eq:v2_usable_window}
\end{equation}
式\eqref{eq:v2_usable_window}不是普适的遮挡阈值，而是依赖 $N$、神经场参数、输入制度、有效轨迹集合、$q_\omega$、$R$、初始协方差和误差容许值的条件函数。它的作用是把算法中的 $\tau_{\max}$ 和 $P_{11,\max}$ 与理论误差尺度联系起来，而不是替代独立的闭环实验。

\subsection{理论结果的验证方式与边界}

本章理论至少需要通过以下三类检查验证。第一，在已保存方位轨迹上比较实测方位误差、$V(T)$ 和 $\gamma_1(T)$，检查理论趋势是否成立。第二，在相同观测历史下计算 KF-mean 与 KF-uncertainty 的社会输入均方误差，检查式\eqref{eq:v2_mean_risk_gap}给出的风险差方向。第三，在完整检查点分叉实验中记录神经场状态差、速度差和位置差，检查式\eqref{eq:v2_field_finite_bound}是否成为保守但有效的上界。

上述理论仍有明确边界：
\begin{enumerate}
 \item wrapped normal 只描述单峰、局部连续的方位后验，不能覆盖任意大方差、多峰后验和长时间周期分支切换；
 \item 条件均方最优性只针对相同有效轨迹集合，不能解决邻居存在性和错误数据关联；
 \item 方向信息窗口只描述周期输入的一阶谐波，不是 GO、LO、凝聚或碰撞安全窗口；
 \item 有限时间神经场界需要局部 Lipschitz 条件，不能直接升级为全局闭环稳定性证明；
 \item 协方差校准不足时，理论 $V(T)$ 仍可能低估真实方位误差；
 \item 位置界只描述同一检查点分叉后的个体轨迹偏差，不能替代群体层面的统计实验。
\end{enumerate}

\subsection{本章理论贡献总结}

本章的理论贡献可以概括为四个层次。第一，在给定 wrapped normal 方位后验和相同有效轨迹集合的条件下，证明协方差感知周期卷积输入是理想社会输入的条件均方最优估计，并量化 KF-mean 忽略协方差造成的额外重建风险。第二，将连续缺失期间的方位方差传播转化为周期方向谐波衰减，定义可计算的一阶方向信息窗口。第三，区分有效轨迹内的预测误差与完整邻居集合相对于有效轨迹集合的输入缺口，给出完整社会输入误差的可计算分解。第四，在局部 Lipschitz 条件下给出社会输入误差到神经场状态、速度解码和有限时间位置偏差的传播界，从而形成一条不把输入重建最优性冒充为闭环稳定性的理论链条。

因此，本章的创新对象不是 Kalman 经典递推本身，而是“仅方位预测后验到异中心环形神经场输入及其有限时间扰动传播的理论接口”。下一章应分别验证：方位误差是否随理论方差增长、协方差感知输入是否降低有效轨迹集合上的重建风险、漏邻居项是否解释 fresh-drop 的输入缺口，以及理论窗口是否能够预测受控转向实验中的相对退化趋势。

\section{仿真实验与结果}

\subsection{实验设置}

仿真采用 $N=20$、$M=100$、$\Delta t=0.3$、$v_s=0.05$、$\beta=1000$ 和 $\sigma_s=0.4$，周期方形区域边长 $L=1000$。闭环补偿实验采用完整检查点分叉：无遮挡控制和全部遮挡分支从同一位置、航向、神经活动、神经元偏好方向、固定邻居身份及预生成观测噪声流开始。记固定循环有向邻居边数为 $n_e$，实验中取 $n_e=15$，不访问仿真距离进行选邻。

\begin{table}[htbp]
 \centering
 \small
 \caption{扩展配对闭环实验参数}
 \label{tab:simulation_setup}
 \begin{tabular}{ll}
 \toprule
 参数 & 设置\\
 \midrule
 个体数、神经元数 & $N=20$，$M=100$\\
 时间步长、速度尺度 & $\Delta t=0.3$，$v_s=0.05$\\
 神经参数 & $\beta=1000$，$\sigma_s=0.4$，$L=1000$\\
 固定邻居边数 & $n_e=15$\\
 隐藏边数 & $7/15$、$11/15$、$13/15$\\
 连续遮挡时长 & $T=\{3,6,10.2,18,30,60\}$\\
 独立随机种子 & $201$--$230$\\
 测量噪声标准差 & $0.01$ rad\\
 常角速度 KF & $q_\omega=10^{-4}$，$R=10^{-4}$\\
 阻尼角速度 KF & $q_\omega=10^{-5}$，$\tau=10$，$R=10^{-4}$\\
 滤波初始化 & $P_0=\operatorname{diag}(R,1)$\\
轨迹截止 & $\tau_{\max}=120$，$P_{11,\max}=1\ \mathrm{rad}^2$\\
评价窗口 & 遮挡期；遮挡后前 100 步；恢复仿真 300 步\\
受控转向设置 & 3 个领航个体；遮挡前 20 步开始转向；转向角速度 $0.025$ rad/模型时间\\
相对退化指标 & 相对无遮挡控制的 GO、LO 和凝聚距离绝对偏差\\
 \bottomrule
 \end{tabular}
\end{table}

表~\ref{tab:simulation_setup} 中，$T$ 为连续遮挡时长，$q_\omega$ 为角加速度噪声谱强度，$\tau$ 为阻尼模型相关时间，$R$ 为方位测量方差，$P_0$ 为初始化协方差，$\tau_{\max}$ 为观测年龄上限，$P_{11,\max}$ 为方位方差轨迹截止阈值。本节区分四类证据：社会输入机制图使用种子 1--100 的独立机制扫描；预测校准诊断使用种子 138--139 进行参数选择、140--141 进行开发评估；无受控转向的扩展闭环实验和受控转向实验均使用种子 201--230，其中受控转向结果为完整的 $30$ 个种子、$2700$ 条配对闭环记录。主流模型比较仅作为不同传感信息条件下的辅助参考，不与仅方位闭环结果合并统计。

比较五种方法：\textit{fresh-drop} 在缺失时删除外部方位输入；\textit{hold-last} 保持最后一次观测方位；\textit{KF-mean} 使用常角速度 Kalman 预测均值；\textit{KF-uncertainty} 使用常角速度 Kalman 均值和协方差卷积；\textit{KF-GM-uncertainty} 使用阻尼角速度 Kalman 模型和协方差卷积。后文分别将 KF-uncertainty 和 KF-GM-uncertainty 简称为 KF-unc 和 KF-GM。所有方法都保留神经场内部状态，因此 fresh-drop 只表示无外部方位轨迹记忆，不表示完全无记忆。每种条件包含 30 个种子，共 2700 条配对闭环结果。

\subsection{社会输入机制的基准影响}

图~\ref{fig:input_mechanism} 同时给出输入幅值、角度覆盖、输入集中度、GO、LO 和凝聚距离。随机选邻并采用原始单邻居幅值时，增加 $n_e$ 会同时提高输入幅值和方向采样密度，使 GO、LO 上升并减小凝聚距离。固定总幅值后，输入幅值基本保持不变，但角度覆盖和输入集中度仍随选邻策略变化；因此减少邻居的主要代价表现为方向信息不足。该结果支持将输入幅值和方位结构作为两个相互耦合但可区分的因素，也说明遮挡补偿不能只通过维持输入质量解释。

\begin{figure}[htbp]
 \centering
 \includegraphics[width=0.96\linewidth]{fig1_input_mechanism.png}
 \caption{邻居数量、输入幅值和方位选择策略的机制分解（N=20，种子 1--100）。}
 \label{fig:input_mechanism}
\end{figure}

\subsection{连续缺失下的协方差与方位预测}

连续缺失期间，常角速度模型的方位协方差按
\begin{equation}
 V(T)=p_{11}+2Tp_{12}+T^2p_{22}+q_\omega T^3/3
\end{equation}
增长。周期卷积对第 $n$ 阶方向谐波的衰减为 $\exp(-n^2V/2)$，因此协方差卷积不会删除预测输入，而是随着缺失时间增加展宽方向核并抑制高阶方向结构。图~\ref{fig:uncertainty_diagnostics} 的上排给出该解析性质，下排给出离线预测误差和 95\% 经验覆盖率。覆盖率在较长遮挡下明显低于名义值，说明当前协方差仍偏乐观；均值预测改善不能替代不确定性校准。该结果也限定了后续闭环解释：协方差卷积的收益应理解为对过期方向输入的抑制，而不是已经校准的概率保证。

\begin{figure}[htbp]
 \centering
 \includegraphics[width=0.98\linewidth]{kf-uncertainty-diagnostics.png}
 \caption{方位预测诊断。上排为解析协方差与方向信息保持；下排为离线均值误差和 95\% 经验覆盖率（种子 138--141）。}
 \label{fig:uncertainty_diagnostics}
\end{figure}

图~\ref{fig:predict_error} 进一步给出隐藏边上的平均方位预测误差。相较 hold-last，阻尼角速度 KF-GM-uncertainty 在 $T\ge 6$ 时将预测误差降低约 $13\%$--$19\%$，配对 bootstrap 区间不跨零；在 $T=60$ 时仍降低约 $9\%$--$10\%$。$T=3$ 时均值降低约 $11\%$--$13\%$，但在两种较强遮挡条件下区间略跨零，因此只视为方向一致的短期改善。常角速度 KF 在短时遮挡中与 hold-last 接近，在中等和长时间缺失中才逐渐表现出均值预测优势。

\begin{figure}[htbp]
 \centering
 \includegraphics[width=0.96\linewidth]{kf-prediction-error.png}
 \caption{隐藏有向边上的方位预测误差。误差仅用于诊断，不进入控制器。}
 \label{fig:predict_error}
\end{figure}

\subsection{闭环补偿结果}

为避免长遮挡下轨迹分叉导致 GO 符号偶然变化，主响应采用相对无遮挡控制的绝对偏差。记 $\mathcal K_o$ 为一次连续遮挡覆盖的时间步集合，遮挡期 GO 偏差定义为
\begin{equation}
 D_{\mathrm{GO}}=\frac{1}{|\mathcal K_o|}\sum_{k\in\mathcal K_o}
 \left|GO_{\mathrm{control}}(k)-GO_{\mathrm{method}}(k)\right|,
\end{equation}
遮挡后凝聚距离偏差采用相同定义。配对改善量由同一随机种子的方法差值与 fresh-drop 损失计算，区间由 5000 次 bootstrap 重采样得到。

表~\ref{tab:kf_go} 和表~\ref{tab:kf_pair} 分别给出隐藏 $11/15$ 条输入时 GO 偏差与遮挡后凝聚距离偏差的配对改善。百分比相对 fresh-drop 计算，同时列出 hold-last，以避免把保留最后一次方位本身带来的收益误归因于 Kalman 外推。表中所有相对 fresh-drop 的 95\% bootstrap 区间均不跨零。

\begin{table}[htbp]
 \centering
 \small
 \caption{隐藏 $11/15$ 条输入时 GO 偏差相对 fresh-drop 的降低比例}
 \label{tab:kf_go}
 \begin{tabular}{lcccc}
 \toprule
 $T$ & hold-last & KF-mean & KF-uncertainty & KF-GM-uncertainty\\
 \midrule
 3.0 & 81.6\% & 86.4\% & 86.2\% & 86.5\%\\
 6.0 & 75.0\% & 82.3\% & 82.3\% & 83.3\%\\
 10.2 & 68.6\% & 77.9\% & 77.0\% & 79.5\%\\
 18.0 & 63.5\% & 71.5\% & 67.7\% & 72.7\%\\
 30.0 & 58.2\% & 63.1\% & 52.6\% & 62.6\%\\
 60.0 & 47.9\% & 46.9\% & 33.7\% & 46.4\%\\
 \bottomrule
 \end{tabular}
\end{table}

\begin{table}[htbp]
 \centering
 \small
 \caption{隐藏 $11/15$ 条输入时遮挡后距离偏差相对 fresh-drop 的降低比例}
 \label{tab:kf_pair}
 \begin{tabular}{lcccc}
 \toprule
 $T$ & hold-last & KF-mean & KF-uncertainty & KF-GM-uncertainty\\
 \midrule
 3.0 & 89.3\% & 89.7\% & 89.8\% & 91.3\%\\
 6.0 & 85.5\% & 89.7\% & 89.4\% & 91.2\%\\
 10.2 & 79.3\% & 89.0\% & 86.5\% & 87.0\%\\
 18.0 & 70.1\% & 88.1\% & 83.4\% & 80.1\%\\
 30.0 & 59.0\% & 85.6\% & 79.5\% & 71.5\%\\
 60.0 & 40.2\% & 59.0\% & 64.0\% & 55.3\%\\
 \bottomrule
 \end{tabular}
\end{table}

hold-last 在短中时遮挡中已经解释了大部分偏差降低。Kalman 的额外价值并非在所有指标和时间尺度上都一致：在表~\ref{tab:kf_go} 的代表性条件下，$T\le30$ 时 KF-mean 和 KF-GM 的 GO 偏差低于 hold-last；$T=60$ 时，两者的均值略高于 hold-last，不能宣称长时间 GO 保持仍优于简单保持。相比之下，KF-uncertainty 对遮挡后距离偏差的改善更稳定，在 $T=60$ 时相对 fresh-drop 降低约 $64\%$，也明显高于 hold-last 的 $40.2\%$。这体现了协方差卷积的主要作用：抑制过期方向输入对神经场峰值的错误强化，而不是在所有指标上无条件优于其他方法。

\begin{figure}[htbp]
 \centering
 \includegraphics[width=0.98\linewidth]{kf-performance-summary.png}
 \caption{不同遮挡强度和时长下的绝对 GO 偏差与遮挡后距离偏差。误差线为种子间近似 95\% 置信区间。}
 \label{fig:kf_performance}
\end{figure}

图~\ref{fig:kf_vs_hold} 直接给出各 Kalman 方法相对 hold-last 的配对差值。负值表示相对 hold-last 降低偏差。该图用于区分“保留旧方位已带来的收益”和“Kalman 外推或协方差卷积额外带来的收益”：短时遮挡中多数 Kalman 增量很小，中时遮挡中 KF-mean 和 KF-GM 对 GO 有可见增益，而 KF-uncertainty 的优势主要体现在长时间遮挡后的凝聚结构。

\begin{figure}[htbp]
 \centering
 \includegraphics[width=0.98\linewidth]{kf-versus-hold.png}
 \caption{Kalman 方法相对 hold-last 的配对差值。横轴为遮挡时长，负值表示偏差低于 hold-last。}
 \label{fig:kf_vs_hold}
\end{figure}

\subsection{遮挡前受控转向与相对退化指标}

仅在已经形成一致直线运动后开始遮挡，会使没有预测器的模型仍能依靠已有航向惯性维持较高 GO，因而不能充分检验预测补偿。为此，本文增加受控转向干预。首先从同一无遮挡检查点分叉；在遮挡开始前 $20$ 步（$6$ 个模型时间单位），仅对编号为 $1$--$3$ 的领航个体施加弱的异中心旋转方向输入，领航方向以 $0.025$ rad/模型时间的角速度缓慢变化。遮挡在第 $201$ 步开始，期间领航个体继续转向；其余个体没有额外转向命令，只能通过邻居方位输入跟随。这样，Kalman 轨迹在遮挡前能够接收连续变化的方位观测并估计方位角速度，而遮挡开始时跟随者尚未完全收敛到新的共同方向。

受控转向实验仍比较 fresh-drop、hold-last、KF-mean、KF-uncertainty 和 KF-GM-uncertainty。每个方法与同一转向任务下的无遮挡控制轨迹配对。除绝对 GO、LO 和凝聚距离外，定义相对退化指标
\begin{equation}
 D_{\mathrm{GO}}^{\mathrm{rel}}=
 \frac{\operatorname{mean}_{k\in\mathcal K_o}
 |GO_{\mathrm{method}}(k)-GO_{\mathrm{control}}(k)|}
 {\operatorname{mean}_{k\in\mathcal K_o}|GO_{\mathrm{control}}(k)|},
\end{equation}
并以相同方式定义 $D_{\mathrm{LO}}^{\mathrm{rel}}$ 和遮挡后凝聚距离相对偏差 $D_{\mathrm{pair}}^{\mathrm{rel}}$。这些指标比较的是观测缺失造成的相对退化，而不是不同模型本来具有的绝对性能差异。

图~\ref{fig:turn_timeline} 给出干预时序。图~\ref{fig:turn_relative} 给出完整 $30$ 个种子在 $11/15$ 条有向边隐藏条件下的统计结果。以 $T=3$ 和 $T=10.2$ 为例，fresh-drop 的相对 GO 退化分别为 $0.99\%$ 和 $2.82\%$，KF-GM-uncertainty 分别为 $0.13\%$ 和 $0.45\%$；对应的遮挡后相对凝聚距离偏差分别为 $0.127\%$ 和 $0.605\%$，KF-GM-uncertainty 分别降至 $0.0085\%$ 和 $0.090\%$。完整矩阵表明，遮挡前转向确实放大了无预测补偿分支的退化，而 Kalman 分支在 GO 和凝聚结构指标上均保持较低偏差；误差线为 $30$ 个种子间的近似 $95\%$ 置信区间。

\begin{figure}[htbp]
 \centering
 \includegraphics[width=0.82\linewidth]{intervention-timeline.png}
 \caption{受控转向实验时序：遮挡前先开始领航转向，随后进入连续遮挡和恢复窗口。}
 \label{fig:turn_timeline}
\end{figure}

\begin{figure}[htbp]
 \centering
 \includegraphics[width=0.98\linewidth]{relative-degradation.png}
 \caption{受控转向下相对无遮挡控制的退化（完整 $30$ 个种子）。误差线为种子间近似 $95\%$ 置信区间。}
 \label{fig:turn_relative}
\end{figure}

\subsection{与主流集群模型的参考比较}

表~\ref{tab:mainstream} 给出 $T=30$、隐藏 $7/15$ 条输入时与 Vicsek、Boids 和 Olfati--Saber 型模型的辅助比较。其中，Vicsek 型模型源于自驱动粒子模型\upcite{Vicsek1995}，其宏观集体运动机制可进一步参见相关综述\upcite{Vicsek2012}；Olfati--Saber 型模型采用连续时间一致性框架和势函数构造 flocking 控制\upcite{OlfatiSaber2006}。主流模型使用邻居航向或相对位置，本文方法只使用方位观测，因此该比较用于说明额外状态信息带来的性能上界和代价，而不是同传感器条件下的算法排序。需要特别指出，这一辅助比较不包含与本文等价的方位预测器；若遮挡发生在已经充分收敛的直线运动阶段，主流模型仍可能依靠剩余可见邻居和自身状态维持较高 GO。因此，本文不使用该图单独证明 Kalman 的遮挡收益，Kalman 的因果效果以受控转向配对实验为主要证据。

\begin{table}[htbp]
 \centering
 \small
 \caption{与主流集群模型的参考比较（种子 101--105，$T=30$）}
 \label{tab:mainstream}
 \begin{tabular}{lrrrr}
 \toprule
 模型 & 遮挡期 GO & GO 损失 & 最小距离$/L$ & 危险接近时间占比\\
 \midrule
 本文 KF-uncertainty & 0.729 & $-0.0014$ & 0.00056 & 0.872\\
 Vicsek（邻居航向） & 1.000 & 0 & 0.00760 & 0.800\\
 Boids（位置+航向） & 0.975 & 0.0038 & 0.00773 & 0.064\\
 Olfati--Saber 型（位置+航向） & 0.911 & 0.0368 & 0.00764 & 0.031\\
 \bottomrule
 \end{tabular}
\end{table}

图~\ref{fig:mainstream} 和表~\ref{tab:mainstream} 仅说明不同感知条件下的性能--代价差异：本文方法在只使用方位观测时仍能维持可用的群体有序运动，并保持较低转向控制量；Boids 和 Olfati--Saber 型基线由于使用相对位置并含有排斥或势函数项，在安全距离方面更有优势。不过，原异中心神经场模型没有显式避撞屏障，本文方法的最近距离和危险接近时间占比并不优于这些基线。因此，本文不能被视为安全控制器，最小距离只能作为风险诊断指标。

图~\ref{fig:mainstream} 和表~\ref{tab:mainstream} 中，记 $d_{\min}$ 为最小个体间距离，$d_{\mathrm{safe}}$ 为危险接近诊断阈值，$R_{\mathrm{dp}}$ 为最小距离不超过 $d_{\mathrm{safe}}$ 的时间步比例。上述诊断量不进入控制器，也不构成碰撞安全保证。

\begin{figure}[htbp]
 \centering
 \includegraphics[width=0.96\linewidth]{mainstream-comparison.png}
 \caption{与主流集群模型的参考比较。危险接近采用仿真距离阈值，不等同于具有明确物理半径的碰撞。}
 \label{fig:mainstream}
\end{figure}

\subsection{结果边界}

上述闭环结果说明 Kalman 补偿在给定固定拓扑和参数下能够显著减缓连续缺失引起的性能退化，但仍存在明确边界。第一，协方差卷积的条件均方最优性依赖于给定后验模型，当前预测器的方差校准尚未完成，名义覆盖率不能等同于真实预测可靠性。第二，阻尼角速度模型参数来自前期离线开发，闭环实验没有在这些新种子上重新调参。第三，控制器没有邻居存在概率和数据关联机制，错误预测可能形成幽灵邻居输入。第四，当前实验采用固定身份和脚本化连续遮挡，尚未覆盖真实几何遮挡、动态邻居集合和周期最近像切换。第五，模型没有显式碰撞屏障，因此补偿收益不能扩展为安全保证。

\section{讨论}

本文方法的有效性具有明确边界。第一，方位—角速度 KF 只适用于局部连续角度分支；长时间观测中断、快速转弯和周期边界最近镜像切换会造成模型失配。第二，协方差卷积降低的是方向结构，不会自动恢复从未观测到的邻居，也不能保证有效交互图保持连通。第三，当前实验中的最近邻（nearest）和方位平衡（balanced）策略使用了仿真距离，只能作为机制对照，不能作为仅方位传感器控制器。第四，本文保留原模型的神经活动合矢量运动规则，没有加入显式碰撞屏障函数，因此最小距离只能作为评价指标，不能宣称具有严格碰撞安全保证。第五，本文先采用从 KF 到神经场的单向接口；神经场意图反向调节 KF 预测噪声的双向耦合可以作为后续工作，但需要独立的闭环推导和仿真验证。

从更广的群体运动建模看，流体耦合等机制会进一步改变群体组织方式\upcite{Mao2025}，本文的方位输入模型尚未覆盖这类耦合效应。参数不确定性下的预测控制可通过灵敏度分析显式约束风险和性能退化\upcite{Belvedere2025}，可作为后续将预测误差纳入安全约束的参考。对于分布式部署，还需考虑恶意节点引起的估计偏差\upcite{An2024}以及多跳通信带来的拓扑和资源约束\upcite{Gao2024}。

\section{结论}

本文提出了一个面向间歇方位观测的异中心神经场 Kalman 预测补偿框架。该框架从原始异中心 flocking 模型出发，先分析有效社会输入的数量、幅值和方向结构，再使用方位—方位角速度 Kalman 滤波器在观测缺失期间预测已跟踪邻居的方位，并将预测协方差通过圆周核卷积传递给神经场。理论分析表明，连续缺失会导致方位协方差按时间累积，圆周方向的高阶结构按 $e^{-n^2V/2}$ 衰减；在给定后验模型下，协方差卷积输入具有条件均方最优的输入重建性质。

本文的结果表明，在固定身份、等权输入和遮挡强度可受控的配对实验中，Kalman 预测补偿相对 fresh-drop 显著降低了连续缺失造成的 GO 偏差和遮挡后凝聚距离偏差，并随着遮挡时间增加形成逐步缩短但仍清晰的收益窗口。KF-GM-uncertainty 进一步提高了隐藏边的方位预测精度。本文不将 Kalman 预测解释为普遍稳定性保证，也不把任何单一邻居数量写成普适阈值；预测协方差校准、邻居存在概率、真实几何遮挡和碰撞约束仍是后续工作。

\begin{thebibliography}{30}
\bibitem{Bai2024} C. Bai, P. Yan, H. Piao, W. Pan, and J. Guo, ``Learning-based multi-UAV flocking control with limited visual field and instinctive repulsion,'' \emph{IEEE Transactions on Cybernetics}, vol. 54, no. 1, pp. 462--475, 2024, doi: 10.1109/TCYB.2023.3246985.
\bibitem{Boldrer2026} M. Boldrer, V. Kratky, V. Walter, and M. Saska, ``Distributed Lloyd-based algorithm for uncertainty-aware multi-robot under-canopy flocking,'' \emph{IEEE Robotics and Automation Letters}, vol. 11, no. 5, pp. 5526--5533, 2026, doi: 10.1109/LRA.2026.3673985.
\bibitem{LiFlocking2024} C. Li, Y. Yang, T.-Y. Huang, and X.-B. Chen, ``An improved flocking control algorithm to solve the effect of individual communication barriers on flocking cohesion in multi-agent systems,'' \emph{Engineering Applications of Artificial Intelligence}, vol. 137, art. 109110, 2024, doi: 10.1016/j.engappai.2024.109110.
\bibitem{Ge2025Coordination} X. Ge, Q.-L. Han, X.-M. Zhang, D. Ding, and B. Ning, ``Distributed coordination control of multi-agent systems under intermittent sampling and communication: a comprehensive survey,'' \emph{Science China Information Sciences}, vol. 68, no. 5, art. 151201, 2025, doi: 10.1007/s11432-024-4355-1.
\bibitem{OlfatiSaber2006} R. Olfati-Saber, ``Flocking for multi-agent dynamic systems: algorithms and theory,'' \emph{IEEE Transactions on Automatic Control}, vol. 51, no. 3, pp. 401--420, 2006, doi: 10.1109/TAC.2005.864190.
\bibitem{Vicsek1995} T. Vicsek, A. Czirók, E. Ben-Jacob, I. Cohen, and O. Shochet, ``Novel type of phase transition in a system of self-driven particles,'' \emph{Physical Review Letters}, vol. 75, no. 6, pp. 1226--1229, 1995, doi: 10.1103/PhysRevLett.75.1226.
\bibitem{Salahshour2025} M. Salahshour and I. D. Couzin, ``Allocentric flocking,'' \emph{Nature Communications}, vol. 16, art. 9051, 2025, doi: 10.1038/s41467-025-64676-5.
\bibitem{Xiourouppa2025} A. H. Xiourouppa, D. Mikhin, M. Humphries, and J. Maclean, ``Theoretical insights for bearings-only tracking in log-polar coordinates,'' \emph{IEEE Transactions on Aerospace and Electronic Systems}, vol. 61, no. 4, pp. 8534--8545, 2025, doi: 10.1109/TAES.2025.3547317.
\bibitem{Zhang1996} K. Zhang, ``Representation of spatial orientation by the intrinsic dynamics of the head-direction cell ensemble: a theory,'' \emph{The Journal of Neuroscience}, vol. 16, no. 6, pp. 2112--2126, 1996, doi: 10.1523/JNEUROSCI.16-06-02112.1996.
\bibitem{Chen2026} Z. Chen, H. Wang, J. Li, G. Li, and S. Wang, ``Egocentric head direction encoding and perception model based on an adaptive ring attractor network,'' \emph{IEEE Transactions on Neural Networks and Learning Systems}, vol. 37, no. 8, pp. 3900--3912, 2026, doi: 10.1109/TNNLS.2026.3656687.
\bibitem{Ao2021} P. Ao, ``Towards predictive neural network dynamical theory,'' \emph{Physics of Life Reviews}, vol. 36, pp. 30--32, 2021, doi: 10.1016/j.plrev.2020.08.006.
\bibitem{Lu2025} Z. Lu, J. B. Julian, G. K. Aguirre, and R. A. Epstein, ``A neural compass in the human brain during naturalistic virtual navigation,'' \emph{The Journal of Neuroscience}, vol. 45, no. 34, art. e1765242025, 2025, doi: 10.1523/JNEUROSCI.1765-24.2025.
\bibitem{Krongauz2024} D. L. Krongauz, A. Ayali, and G. A. Kaminka, ``Vision-based collective motion: a locust-inspired reductionist model,'' \emph{PLOS Computational Biology}, vol. 20, no. 1, art. e1011796, 2024, doi: 10.1371/journal.pcbi.1011796.
\bibitem{Li2024} K. Li, L. Li, R. Gross, and S. Zhao, ``A collective perception model for neighbor selection in groups based on visual attention mechanisms,'' \emph{New Journal of Physics}, vol. 26, no. 1, art. 012001, 2024, doi: 10.1088/1367-2630/ad1b81.
\bibitem{Xue2025} K. Xue, D. Kong, P. Wang, Z. Xu, and Z. Huang, ``Attention-based selective interaction model for self-organized collective motion,'' \emph{Chaos, Solitons \& Fractals}, vol. 201, art. 117320, 2025, doi: 10.1016/j.chaos.2025.117320.
\bibitem{Qi2024} J. Qi, Z. Zhang, J. Yu, X. Dong, Q. Li, and H. Jiang, ``Time-varying group formation tracking control for multi-agent systems using distributed multi-sensor multi-target filtering with intermittent observations,'' \emph{International Journal of Robust and Nonlinear Control}, vol. 34, no. 17, pp. 11681--11704, 2024, doi: 10.1002/rnc.7587.
\bibitem{Pan2026} Y. Pan, T. Chen, Q. Qi, H. Tian, and W. Kong, ``Bearing-only target tracking for multi-UUV via IMM-UIF with initial filter value correction,'' \emph{Ocean Engineering}, vol. 356, art. 125347, 2026, doi: 10.1016/j.oceaneng.2026.125347.
\bibitem{Ge2025} Q. Ge, Y. Li, Y. Wang, X. Hu, H. Li, and C. Sun, ``Adaptive Kalman filtering based on model parameter ratios,'' \emph{IEEE Transactions on Automatic Control}, vol. 69, no. 9, pp. 6230--6237, 2024, doi: 10.1109/TAC.2024.3376306.
\bibitem{ZhangBearing2025} K. Zhang, H. Wang, W. Chen, J. Ren, N. Luo, and Z. Lu, ``Multiplatform bearing-only target tracking algorithm based on maximum likelihood maximum correntropy Kalman filter,'' \emph{IEEE Sensors Journal}, vol. 25, no. 7, pp. 11366--11377, 2025, doi: 10.1109/JSEN.2025.3529771.
\bibitem{Yu2024} Q. Yu, Z. Fan, X. Bu, and Z. Hou, ``Event-triggered based predictive iterative learning control with random packet loss compensation for nonlinear networked systems,'' \emph{ISA Transactions}, vol. 148, pp. 169--181, 2024, doi: 10.1016/j.isatra.2024.03.001.
\bibitem{Selvi2025} D. Selvi and G. Battistelli, ``Distributed Kalman filtering with adaptive communication,'' \emph{IEEE Control Systems Letters}, vol. 9, pp. 15--20, 2025, doi: 10.1109/LCSYS.2025.3550401.
\bibitem{Zhou2025} H. Zhou and S. Sun, ``Distributed extended Kalman consensus filtering for multi-sensor networked nonlinear systems under stochastic communication protocol,'' \emph{IEEE Transactions on Signal Processing}, vol. 73, pp. 3629--3640, 2025, doi: 10.1109/TSP.2025.3599848.
\bibitem{Yang2025} T. Yang, J. Qian, Z. Duan, and Z. Sun, ``Distributed Kalman filter with ultimately accurate fused measurement covariance,'' \emph{IEEE Transactions on Automatic Control}, vol. 70, no. 10, pp. 6449--6464, 2025, doi: 10.1109/TAC.2025.3557365.
\bibitem{Rong2024} D. Rong and Y. Wang, ``An adaptive spatial target tracking method based on unscented Kalman filter,'' \emph{Sensors}, vol. 24, no. 18, art. 6094, 2024, doi: 10.3390/s24186094.
\bibitem{Sinopoli2004} B. Sinopoli, L. Schenato, M. Franceschetti, K. Poolla, M. I. Jordan, and S. S. Sastry, ``Kalman filtering with intermittent observations,'' \emph{IEEE Transactions on Automatic Control}, vol. 49, no. 9, pp. 1453--1464, 2004, doi: 10.1109/TAC.2004.834121.
\bibitem{Vicsek2012} T. Vicsek and A. Zafeiris, ``Collective motion,'' \emph{Physics Reports}, vol. 517, no. 3--4, pp. 71--140, 2012, doi: 10.1016/j.physrep.2012.03.004.
\bibitem{Mao2025} X. Mao, J. Wang, X. Mao, and J. Deng, ``A hybrid numerical model for the collective motion of fish groups,'' \emph{Journal of Fluid Mechanics}, vol. 1009, art. A26, 2025, doi: 10.1017/jfm.2025.216.
\bibitem{Belvedere2025} T. Belvedere, M. Cognetti, G. Oriolo, and P. R. Giordano, ``Sensitivity-aware model predictive control for robots with parametric uncertainty,'' \emph{IEEE Transactions on Robotics}, vol. 41, pp. 3039--3058, 2025, doi: 10.1109/TRO.2025.3554415.
\bibitem{An2024} L. An and G.-H. Yang, ``Mean-square exponential convergence for Byzantine-resilient distributed state estimation,'' \emph{Automatica}, vol. 163, art. 111592, 2024, doi: 10.1016/j.automatica.2024.111592.
\bibitem{Gao2024} R. Gao, G.-H. Yang, and S. Wasly, ``Resilient distributed state estimation with multi-hop communication,'' \emph{Automatica}, vol. 168, art. 111823, 2024, doi: 10.1016/j.automatica.2024.111823.
\end{thebibliography}
\end{document}
